\section{Experiments}
\label{sec:experiments}

The experimental arc proceeds in five stages.  Predictive deep
supervision makes the looped encoder trainable at intermediate depths.
A discriminability probe shows that high lift is not a sufficient
success signal because the final readout can collapse.  VICReg at the
readout resolves that collapse while preserving lift.  A
full-validation audit corrects a $12\times$ inflation in the
small-gallery KNN metric and surfaces a peak-then-overfit curve.
Finally, a CPU retrieval head-to-head places the $1.3$M \CodeWM{}
within noise of $125.9$M UnixCoder and well ahead of $124.6$M
CodeBERT.

\subsection{Deep Supervision Raises Lift}
\label{subsec:deepsup}

A short 400-step pass across 15 candidate recipes isolates two
interventions that consistently improve prediction quality: predictive
auxiliary supervision at intermediate loops, and reducing the encoder
to three loops.  Contrastive auxiliary losses and dense skip connections
do not help at this scale.

\begin{table}[t]
\centering
\caption{Selected 400-step screening results.  Predictive aux at loops
$\{1,2,3\}$ beats the \CodeWM{} control by a wide margin on delta
cosine.}
\label{tab:wave1}
\smallskip
\begin{tabular}{llll}
\toprule
Run & Loss & Delta cosine & Category \\
\midrule
\KernelWM{} baseline & $0.017$ & $0.973$ & Kernel control \\
aux-pred & $0.024$ & $0.955$ & Predictive aux at loops 1,2,3 \\
loops3 & $0.051$ & $0.912$ & Shallow encoder \\
\KernelWM{} + aux & $0.052$ & $0.914$ & Kernel aux variant \\
\CodeWM{} baseline & $0.063$ & $0.889$ & Code control \\
\bottomrule
\end{tabular}
\end{table}

Extending the five best recipes to 2K training steps, the combination
of a three-loop encoder with predictive auxiliary supervision
(\code{loops3-auxpred}) is the clear winner on lift over the copy
baseline:

\begin{table}[t]
\centering
\caption{Promoted variants at 2K steps.  Lift is prediction cosine
minus copy-baseline cosine.}
\label{tab:wave2}
\smallskip
\begin{tabular}{lllll}
\toprule
Run & Loss & Delta cosine & Lift & Copy cosine \\
\midrule
\KernelWM{} baseline & $0.003$ & $0.997$ & $+0.051$ & $0.949$ \\
loops3-auxpred & $0.008$ & $0.987$ & $+0.686$ & $0.295$ \\
aux-pred & $0.013$ & $0.979$ & $+0.638$ & $0.336$ \\
loops3 & $0.047$ & $0.903$ & $+0.086$ & $0.895$ \\
aux-loop1 & $0.050$ & $0.975$ & $+0.353$ & $0.617$ \\
\bottomrule
\end{tabular}
\end{table}

On the initial seed (42, pre-fix), the recipe appeared to scale with
training budget: lift crossed $+1.0$ by step 8100.  A post-hoc
investigation (\S\ref{subsec:cross_seed}) found that earlier
multi-seed lift agreement ($\approx +0.82$ at step 6K across seeds
42/43/44) was measured on a batch-128 metric that does not track
full-validation retrieval, and the peak-seed run benefited from a
fortunate initialization due to a seed bug.  The lift trajectory
nonetheless motivated the deep-supervision recipe, even though
the cross-seed retrieval story is weaker than initially assumed.

\subsection{Lift Alone Is a Gameable Metric}
\label{subsec:gameable}

The follow-up discriminability probe changed the interpretation.  On 100
real validation samples, the high-lift \code{loops3-auxpred} checkpoint
had an online effective rank of only $1.95/128$ and KNN@1 of $0.04$.
The predictor had high diagonal alignment with the target, but also high
off-diagonal alignment with wrong targets.

\begin{table}[t]
\centering
\caption{Readout-collapse diagnostic on the high-lift
\code{loops3-auxpred} checkpoint.}
\label{tab:readout_collapse}
\smallskip
\begin{tabular}{lll}
\toprule
Metric & Value & Interpretation \\
\midrule
Online effective rank & $1.95/128$ & Final readout is nearly 2D \\
Target effective rank & $8.27/128$ & Target is only partially healthier \\
Predicted effective rank & $2.77/128$ & Predictions are heavily clustered \\
Cross diagonal cosine & $0.97$ & Correct pairs look aligned \\
Cross off-diagonal cosine & $0.89$ & Wrong pairs also look aligned \\
KNN@1 over 100 examples & $0.04$ & Barely above random \\
\bottomrule
\end{tabular}
\end{table}

This does not erase the deep-supervision result.  Per-loop probes showed
that the encoder intermediates still preserved edit information: loop 3
edit classification reached $94.5\%$, and final delta/state ratio
improved from the pre-supervision value $0.010$ to about $0.033$.  The problem
was narrower and more specific: the attention-pool readout compressed a
rich token representation into a low-rank embedding.

\subsection{VICReg Repairs the Readout}
\label{subsec:vicreg}

More than 30 anti-collapse variants were tested after the caveat:
dense skip connections, dropout, stronger covariance penalties, and
early fusion alone all failed to restore discriminability.  The
intervention that worked is VICReg at the readout combined with the
predictive auxiliary supervision from \S\ref{subsec:deepsup}.

\begin{table}[t]
\centering
\caption{VICReg variants at 2K steps, sorted by online effective
rank.  KNN values use a 128-sample gallery (random KNN@1 baseline
$3.9\%$); corrected full-validation numbers appear in
\S\ref{subsec:fullval}.}
\label{tab:vicreg_screen}
\smallskip
\begin{tabular}{llllll}
\toprule
Run & Lift & KNN@1 & KNN@5 & KNN@10 & Eff. rank \\
\midrule
w7-vicreg-s43 & $+1.16$ & $0.086$ & $0.242$ & $0.422$ & $53.80$ \\
w5b-vicreg-pred & $+1.06$ & $0.086$ & $0.359$ & $0.477$ & $53.52$ \\
w6-vicreg-pred-base & $+1.12$ & $0.086$ & $0.266$ & $0.391$ & $52.12$ \\
w6-vicreg-loops6 & $+1.05$ & $0.039$ & $0.234$ & $0.477$ & $51.66$ \\
w6-vicreg-contrast-aux & $+1.22$ & $0.141$ & $0.320$ & $0.500$ & $51.44$ \\
\bottomrule
\end{tabular}
\end{table}

The strongest all-around variant is \code{w5b-vicreg-pred}: it
preserves lift at $+1.06$ while moving online effective rank from
$\sim\!2$ to $53/128$ and KNN@5 from $0.04$ to $0.36$.  These numbers
guided recipe selection, but the 128-sample gallery baseline ($3.9\%$)
is far above the full-validation baseline ($0.02\%$).
\S\ref{subsec:fullval} reports the corrected numbers.

\subsection{Full-Validation Audit and Overfit Curve}
\label{subsec:fullval}

The batch-128 KNN numbers in Table~\ref{tab:vicreg_screen} use a 128-sample
gallery, so the random KNN@1 baseline is $5/128 \approx 3.9\%$.  Moving to
the full validation set ($N=5000$) shifts the random baseline down by
roughly $40\times$ to $0.02\%$ and reveals both a tighter VICReg-versus-baseline
contrast and a clear overfit beyond step 28K.

\begin{table}[t]
\centering
\caption{Full-validation audit (N=5000, random KNN@1 baseline $0.02\%$).
Pre-VICReg \code{loops3} baseline contrasted with six \code{vicreg\_promotion}
checkpoints from a single seed-42 run.  Row marked $\star$ is the peak.}
\label{tab:fullval_curve}
\smallskip
\resizebox{\textwidth}{!}{%
\begin{tabular}{llllllllll}
\toprule
Checkpoint & Steps & VICReg & KNN@1 & KNN@5 & KNN@10 & KNN@50 & eff.~rank & gap \\
\midrule
\code{loops3} baseline & 5.5K & off & $0.0008$ & $0.0040$ & $0.0096$ & $0.0364$ & $2.09$ & $0.088$ \\
\code{vicreg\_sota} & 14.5K & on & $0.0106$ & $0.0342$ & $0.0592$ & $0.1810$ & $55.79$ & $0.175$ \\
\code{vicreg\_promo} & 20K & on & $0.0092$ & $0.0422$ & $0.0660$ & $0.1710$ & $61.54$ & $0.126$ \\
\code{vicreg\_promo} $\star$ & 28K & on & $0.0142$ & $0.0580$ & $0.0990$ & $0.2608$ & $53.29$ & $0.120$ \\
\code{vicreg\_promo} & 36K & on & $0.0120$ & $0.0380$ & $0.0678$ & $0.1972$ & $53.43$ & $0.106$ \\
\code{vicreg\_promo} & 42K & on & $0.0096$ & $0.0428$ & $0.0818$ & $0.2694$ & $51.33$ & $0.105$ \\
\code{vicreg\_promo} & 44K & on & $0.0026$ & $0.0224$ & $0.0468$ & $0.1978$ & $51.83$ & $0.103$ \\
\bottomrule
\end{tabular}%
}
\end{table}

\begin{figure}[t]
\centering
% Scaling curve for tab:fullval_curve — PGFPlots.
% Data: full-val N=5000 on seed-42 loops3+aux_pred+VICReg+EMA 0.99999.
\begin{tikzpicture}
\begin{axis}[
    name=leftplot,
    width=0.5\textwidth,
    height=4.4cm,
    xlabel={Training step (K)},
    ylabel={Full-val KNN@$k$},
    xmin=0, xmax=48,
    ymin=0, ymax=0.11,
    ytick={0,0.02,0.04,0.06,0.08,0.10},
    xtick={5.5,14.5,20,28,36,42,44},
    xticklabel style={font=\scriptsize},
    yticklabel style={font=\scriptsize},
    xlabel style={font=\small},
    ylabel style={font=\small},
    grid=major,
    grid style={dashed, gray!25},
    legend pos=north east,
    legend style={font=\scriptsize, fill=white, fill opacity=0.85, draw=none, inner sep=2pt},
    legend cell align={left},
    axis line style={-},
    tick align=outside,
    enlargelimits=false,
    clip marker paths=true,
]
\addplot+[mark=o, thick, color=blue!80!black]
    coordinates {(5.5,0.0008) (14.5,0.0106) (20,0.0092) (28,0.0142) (36,0.0120) (42,0.0096) (44,0.0026)};
\addlegendentry{KNN@1}
\addplot+[mark=square*, thick, color=green!50!black]
    coordinates {(5.5,0.0040) (14.5,0.0342) (20,0.0422) (28,0.0580) (36,0.0380) (42,0.0428) (44,0.0224)};
\addlegendentry{KNN@5}
\addplot+[mark=triangle*, thick, color=red!80!black]
    coordinates {(5.5,0.0096) (14.5,0.0592) (20,0.0660) (28,0.0990) (36,0.0678) (42,0.0818) (44,0.0468)};
\addlegendentry{KNN@10}
\draw[dashed, gray!60] (axis cs:28,0) -- (axis cs:28,0.11);
\node[gray!70!black, anchor=south west, font=\scriptsize] at (axis cs:28.5,0.095) {peak};
\draw[dotted, gray!60] (axis cs:5.5,0) -- (axis cs:5.5,0.11);
\node[gray!70!black, anchor=south west, font=\scriptsize] at (axis cs:5.8,0.095) {pre-VICReg};
\end{axis}

\begin{axis}[
    at={($(leftplot.east)+(1.1cm,0)$)}, anchor=west,
    width=0.5\textwidth,
    height=4.4cm,
    xlabel={Training step (K)},
    ylabel={Online effective rank},
    xmin=0, xmax=48,
    ymin=0, ymax=130,
    ytick={0,32,64,96,128},
    xtick={5.5,14.5,20,28,36,42,44},
    xticklabel style={font=\scriptsize},
    yticklabel style={font=\scriptsize},
    xlabel style={font=\small},
    ylabel style={font=\small},
    grid=major,
    grid style={dashed, gray!25},
    axis line style={-},
    tick align=outside,
    enlargelimits=false,
]
\addplot+[mark=o, thick, color=violet]
    coordinates {(5.5,2.09) (14.5,55.79) (20,61.54) (28,53.29) (36,53.43) (42,51.33) (44,51.83)};
\draw[dashed, gray!60] (axis cs:28,0) -- (axis cs:28,130);
\draw[dotted, gray!60] (axis cs:5.5,0) -- (axis cs:5.5,130);
\draw[gray!50, thin] (axis cs:0,128) -- (axis cs:48,128);
\node[gray!70!black, anchor=south east, font=\scriptsize] at (axis cs:46,121) {dim = 128};
\end{axis}
\end{tikzpicture}

\caption{Full-validation scaling curve for the \code{loops3 + aux\_pred +
VICReg} recipe (seed 42, N=5000 gallery).  \textbf{Left}: retrieval
KNN@$\{1,5,10\}$ peaks sharply at step 28K, then degrades monotonically as
the online encoder overfits.  \textbf{Right}: online effective rank jumps
from $2.09$ (pre-VICReg, step 5.5K) to $55.79$ by step 14.5K, peaks at
$61.54$ around step 20K, and then compresses back.  Roughly 80 epochs
over the 45K training split is the ceiling for this recipe.}
\label{fig:scaling_curve}
\end{figure}

\paragraph{Seed-bug caveat.}
The scaling curve above was produced from a single seed-42 run that,
post-hoc, was found to benefit from a fortunate initialization caused by
a seed bug in the training harness.  The data in
Table~\ref{tab:fullval_curve} are correct measurements of that
checkpoint.  Cross-seed results that contextualize the peak appear in
\S\ref{subsec:cross_seed}.

Three findings follow.  First, VICReg is a real readout fix: KNN@5
improves by $14\times$ and effective rank by $25\times$ at the step-28K
peak relative to the pre-VICReg baseline --- a $\sim\!20\times$
headline improvement in readout discriminability, stronger on full-val
than on the small-gallery probe.  Second, the $12\times$ KNN@5 inflation between
batch-128 and full-val is entirely a gallery-size artifact: the
ranking itself is unchanged, but a 128-sample random baseline of
$3.9\%$ made every absolute number look $\sim\!12\times$ better than
it is against the $0.02\%$ full-val baseline.  Third, online-encoder
KNN peaks at step 28K and degrades monotonically afterward (KNN@1
drops from $0.0142$ at 28K to $0.0026$ at 44K) even as the EMA
target's effective rank continues to climb.

\paragraph{Data-to-capacity ratio as the binding constraint.}
Peak at step 28K corresponds to roughly $80$ epochs over the $45$K
CommitPackFT~\citep{muennighoff2024octopack} training split.  Beyond
that point the model memorizes
and held-out discriminability collapses.  The EMA target is unaffected
because it tracks the online encoder slowly and benefits from the
averaging.  This suggests that for tiny JEPA at this recipe, the
binding constraint is data volume relative to model capacity, not
architecture depth: doubling the training corpus is a more promising
axis than scaling the encoder.

\subsection{External Benchmark vs Code Foundation Models}
\label{subsec:foundation}

To place the recipe on external ground, we compare the peak step-28K
\CodeWM{} checkpoint against UnixCoder-base~\citep{guo2022unixcoder}
and CodeBERT-base~\citep{feng2020codebert} --- both of which were
originally evaluated on code-to-code retrieval in the
CodeSearchNet~\citep{husain2019codesearchnet} tradition --- on
CommitPackFT Python edit-pair retrieval at two scales: a
\emph{development-scale} pass (100 queries $\times$ 300 gallery,
step-9.5K \CodeWM{} checkpoint) and a \emph{paper-grade} pass (1000
queries $\times$ 5000 gallery, step-28K \CodeWM{}).  Retrieval uses the
action-delta cosine for \CodeWM{} (target\_encoder of after, minus
online\_encoder of before) and mean-pool of \code{last\_hidden\_state}
deltas for the foundation models, the standard HuggingFace
sentence-embedding recipe.  Timings are single-threaded PyTorch CPU
inference at batch size $1$ on tokenized inputs truncated to the
respective model maxima ($512$ for both foundation models; $\leq 512$
tokens per side for \CodeWM{}), averaged over the query commits on
the same hardware for all three models.

\begin{table}[t]
\centering
\caption{Development-scale retrieval head-to-head on CommitPackFT
Python (100 queries $\times$ 300 gallery, seed 42, \CodeWM{} at
step 9.5K).}
\label{tab:foundation_bench}
\smallskip
\begin{tabular}{lllllll}
\toprule
Model & Params & CPU ms & MRR joint & MRR cos$_{0.95}$ & R@1 joint & R@1 cos$_{0.95}$ \\
\midrule
\CodeWM{} (ours) & $1.3$M & $8.4$ & $\mathbf{0.850}$ & $\mathbf{0.758}$ & $\mathbf{0.76}$ & $\mathbf{0.63}$ \\
UnixCoder-base & $125.9$M & $83.7$ & $0.838$ & $0.724$ & $0.75$ & $0.58$ \\
CodeBERT-base & $124.6$M & $175.8$ & $0.724$ & $0.652$ & $0.57$ & $0.50$ \\
\bottomrule
\end{tabular}
\end{table}

\begin{table}[t]
\centering
\caption{Paper-grade retrieval head-to-head on CommitPackFT Python
(1000 queries $\times$ 5000 gallery, seed 42, \CodeWM{} at the peak
step-28K checkpoint).  \code{by\_edit\_type} saturates to $\approx\!1.0$
for both models and is omitted.  CodeBERT was not rerun at this scale;
development-scale results above suggest the gap would widen against it.}
\label{tab:foundation_bench_paper}
\smallskip
\resizebox{\textwidth}{!}{%
\begin{tabular}{llllllll}
\toprule
Model & Params & MRR joint & MRR cos$_{0.9}$ & MRR cos$_{0.95}$ & R@1 joint & R@1 cos$_{0.9}$ & R@1 cos$_{0.95}$ \\
\midrule
\CodeWM{} (ours) & $1.3$M & $\mathbf{0.790}$ & $\mathbf{0.742}$ & $0.662$ & $\mathbf{0.668}$ & $\mathbf{0.595}$ & $0.486$ \\
UnixCoder-base & $125.9$M & $0.760$ & $0.736$ & $\mathbf{0.671}$ & $0.609$ & $0.574$ & $\mathbf{0.490}$ \\
\bottomrule
\end{tabular}%
}
\end{table}

Scaling the benchmark gallery by $\sim\!16\times$ shrinks the gap.  On
\code{by\_joint}, \CodeWM{} keeps a $+0.030$ MRR and $+0.059$ R@1 edge
over UnixCoder, and the efficiency advantage is unchanged ($100\times$
smaller, $10\times$ lower CPU latency).  On the moderate action-cosine
threshold \code{cos@0.9} the edge is within noise ($+0.006$ MRR,
$+0.021$ R@1).  On the \emph{strictest} threshold \code{cos@0.95},
UnixCoder wins by $-0.009$ MRR and $-0.004$ R@1, a statistical tie in
\CodeWM{}'s favor at development scale that inverts in UnixCoder's
favor at paper scale.  The honest summary is that \CodeWM{} matches
UnixCoder on paper-grade edit retrieval --- moderate wins on moderate
criteria, tie-to-slight-loss on the strictest --- while delivering the
same task at two orders of magnitude less capacity and an order of
magnitude less CPU latency.  Cross-seed variance on the underlying recipe is reported in
\S\ref{subsec:cross_seed}.  A CodeT5+~\citep{wang2023codet5plus}
retry via its dedicated embedding head remains a pending check.

\subsection{Cross-Seed Variance and Predictor Collapse}
\label{subsec:cross_seed}

A seed bug in the training harness was identified post-hoc: the
reported SOTA checkpoint (seed~42, pre-fix, step~28K) benefited from a
fortunate random initialization.  After fixing the bug, the same recipe
was rerun under seeds 42 and 43.  Infrastructure drift was ruled out:
git SHA, config, data pipeline, and early loss curves all match between
pre-fix and post-fix runs.  Table~\ref{tab:cross_seed} reports the
full-validation matrix.

\begin{table}[t]
\centering
\caption{Cross-seed full-validation matrix (N=5000).  The pre-fix
seed-42 peak ($\star$) is the checkpoint used throughout the paper.
Post-fix runs trail by $2\text{--}3\times$ on KNN@5.  Predictor
effective rank is uniformly low ($5\text{--}6/128$) across all runs,
including the peak.  Random baseline: KNN@1=$0.02\%$, KNN@5=$0.10\%$.}
\label{tab:cross_seed}
\smallskip
\resizebox{\textwidth}{!}{%
\begin{tabular}{llllllllll}
\toprule
Run & Step & Seed & Bug status & Eff.\ rank (enc) & Eff.\ rank (pred) & KNN@1 & KNN@5 & KNN@10 & KNN@50 \\
\midrule
\code{vicreg\_promotion} $\star$ & 28K & 42 & pre-fix & $53.3$ & $6.1$ & $1.42\%$ & $5.80\%$ & $9.90\%$ & $26.1\%$ \\
\code{vicreg\_s43} & 28K & 43 & pre-fix & $58.7$ & $5.9$ & $0.50\%$ & $2.20\%$ & $3.98\%$ & $15.4\%$ \\
\code{vicreg\_fixed\_s42} & 16K & 42 & post-fix & $59.9$ & $5.2$ & $0.82\%$ & $2.42\%$ & $4.64\%$ & $18.6\%$ \\
\code{vicreg\_fixed\_s43} & 6K & 43 & post-fix & $63.6$ & $5.9$ & $0.34\%$ & $1.48\%$ & $2.98\%$ & $9.98\%$ \\
Random baseline & --- & --- & --- & --- & --- & $0.02\%$ & $0.10\%$ & $0.20\%$ & $1.00\%$ \\
\bottomrule
\end{tabular}%
}
\end{table}

Three findings emerge.  First, the peak seed-42 pre-fix result
(KNN@5=$5.80\%$) is a genuine outlier: the next-best seed trails at
$2.42\%$, a $2.4\times$ gap.  The recipe consistently beats the random
baseline by $15\text{--}58\times$ on KNN@5, but the headline number
should be understood as a single-seed peak with substantial
cross-seed variance.  Second, encoder effective rank is healthy across
all runs ($53\text{--}64/128$), confirming that VICReg repairs the
encoder readout regardless of seed.  Third, and most important,
predictor effective rank is uniformly collapsed at
$5\text{--}6/128$ across \emph{every} checkpoint, including the peak.
The predictor is the bottleneck.  Encoder representations are
high-rank and discriminative; the predictor compresses them into a
low-dimensional prediction subspace.  This was hidden by lift-only
metrics and by the single-seed focus of the original evaluation.

\paragraph{Per-loop edit-type probe.}
A linear probe for edit-type classification (3 classes) on each loop's
pooled output confirms that the encoder preserves shallow features
throughout:

\begin{center}
\smallskip
\begin{tabular}{lcccccc}
\toprule
Checkpoint & Loop 0 & Loop 1 & Loop 2 & Final & Target & Straightness \\
\midrule
s42 pre-fix 28K $\star$ & $94.5$ & $94.5$ & $99.25$ & $98.5$ & $67.75$ & $0.213$ \\
s42 post-fix 16K & $96.0$ & $94.25$ & $96.5$ & $94.75$ & $69.5$ & $0.368$ \\
\bottomrule
\end{tabular}
\smallskip
\end{center}

Edit type is already recoverable from loop~0 ($92\text{--}96\%$),
so it is a ``shallow'' feature.  The straightness score measures how
linearly the latent trajectory progresses across loops: the post-fix
run is \emph{straighter} ($0.368$ vs.\ $0.213$) despite performing
worse on KNN.  This counterintuitive direction suggests that raw
trajectory straightness is not the binding factor for downstream
retrieval at this scale.  The connection to the temporal straightening
literature~\citep{straightening2026}, where curvature penalties promote
linear latent paths, remains suggestive but requires more loops
(the current 3-loop encoder yields only a single triple) to be
diagnostic.

\subsection{Remaining Checks}

The load-bearing remaining checks are:
\begin{enumerate}
    \item retry CodeT5+~\citep{wang2023codet5plus} with its dedicated
          embedding-head API;
    \item test whether the recipe transfers to \KernelWM{} without
          memorizing the smaller kernel corpus;
    \item investigate predictor-specific regularization to raise
          predictor effective rank above the current $5\text{--}6/128$
          ceiling.
\end{enumerate}
